% !TEX program = lualatex
% ---------------------------------------------------------------------------------------------
%  Deformalization 2: local nilpotence over a cyclotomic tower   (riemann-rev)
%
%  Build:  python3 -B scripts/build_digest.py local-nilpotence
%
%  One part of a unit of the Lean archive behind L. Alpöge's announced proof of RH, rewritten as ordinary
%  mathematics. The prose is hand-written. Lean names, line ranges and counts come from data.tex; the quoted
%  Lean statements are copied from the archive by the script. The note was written from the reading project's
%  structured proofs; the Lean files of the unit were not opened (§ "Notes on the formal text"). For the edition
%  under the signature rule the definitions were read in the Lean source, the proofs were not.
%
%  Rules of the text: standard terminology; every symbol is explained where it first appears and listed again
%  in the notation at the end; no letter has two meanings. The signature rule: a function that the archive
%  defines is written with every explicit argument of its Lean definition, in every place, through the macros
%  that the build writes into data.tex from digests/notation.json and accounts.json (the text never defines
%  them); every hypothesis of a theorem stands inside its statement. The archive's closed normal subgroup U is
%  H here, and its U in the module of characters is C.
%
%  Layout: the reading edition of the guide (11 pt Libertinus on a 140 mm measure); preamble as in the first note.
% ---------------------------------------------------------------------------------------------
\documentclass[11pt,a4paper]{article}
\ifdefined\pdfvariable\pdfvariable suppressoptionalinfo \numexpr 1+2+4+8\relax\fi
\usepackage[a4paper,hmargin=35mm,top=22mm,bottom=27mm,footskip=12mm,nohead]{geometry}

\usepackage{fontspec}
\directlua{
  luaotfload.add_fallback("gfallback", {"STIXTwoMath-Regular.otf:mode=node;", "DejaVuSans.ttf:mode=node;"})
  luaotfload.add_fallback("gmonofallback", {"DejaVuSansMono.ttf:mode=node;", "STIXTwoMath-Regular.otf:mode=node;"})
}
\setmainfont{LibertinusSerif}[Extension=.otf,UprightFont=*-Regular,ItalicFont=*-Italic,
  BoldFont=*-Semibold,BoldItalicFont=*-SemiboldItalic,RawFeature={fallback=gfallback}]
\setsansfont{Inter}[Extension=.otf,UprightFont=*-Regular,ItalicFont=*-Italic,
  BoldFont=*-SemiBold,BoldItalicFont=*-SemiBoldItalic,RawFeature={fallback=gfallback}]
\setmonofont{JetBrainsMono}[Extension=.otf,UprightFont=*-Regular,ItalicFont=*-Italic,
  BoldFont=*-SemiBold,BoldItalicFont=*-SemiBoldItalic,Scale=MatchLowercase,
  RawFeature={-calt;-liga;fallback=gmonofallback}]
\usepackage{amsmath}
\usepackage{amsthm}
\usepackage[math-style=TeX]{unicode-math}
\setmathfont{LibertinusMath-Regular.otf}

\usepackage{xcolor}
\usepackage{tikz}
\usetikzlibrary{calc,arrows.meta,positioning,backgrounds}
\usepackage[protrusion=true,expansion=true]{microtype}
\usepackage{enumitem}
\usepackage{needspace}
\usepackage{booktabs}
\usepackage{array}
\usepackage{longtable}
\usepackage{fancyhdr}
\usepackage{lastpage}
\usepackage{multicol}
\usepackage[hidelinks,pdfusetitle]{hyperref}

\definecolor{ink}{HTML}{0F172A}
\definecolor{inksoft}{HTML}{475569}
\definecolor{inkfaint}{HTML}{94A3B8}
\definecolor{rulegrey}{HTML}{CBD5E1}
\definecolor{arith}{HTML}{C2410C}
\definecolor{est}{HTML}{7C3AED}
\definecolor{ana}{HTML}{0891B2}
\definecolor{ok}{HTML}{047857}
\definecolor{leancol}{HTML}{3730A3}

\title{Local nilpotence over a cyclotomic tower}
\author{riemann-rev}
\date{}

% ------------------------------------------------------------------------------------- macros
\newcommand{\leanfont}{\ttfamily\fontsize{8.6}{10.5}\selectfont\color{leancol}}
\newcommand{\lean}[1]{\csname LN@#1\endcsname}
\newcommand{\leanquote}[1]{\csname LQ@#1\endcsname}
% generated by scripts/build_digest.py -- do not edit
\expandafter\newcommand\csname LN@XDual\endcsname{\href{https://github.com/Deicyde/riemann-hypothesis/blob/40e56bab026552e089d7afcfeff08d2defb9c005/src/Development/AutomorphicArithmetic/AutomorphicVarieties/GaloisLibrary/KummerTheory/Mk0IdMapCongr.lean\#L90-L99}{\leanfont Kum.XDual}}
\expandafter\newcommand\csname LN@Dinf\endcsname{\href{https://github.com/Deicyde/riemann-hypothesis/blob/40e56bab026552e089d7afcfeff08d2defb9c005/src/Development/AutomorphicArithmetic/AutomorphicVarieties/GaloisLibrary/KummerTheory/Mk0IdMapCongr.lean\#L195-L195}{\leanfont Kum.Dinf}}
\expandafter\newcommand\csname LN@Xinf\endcsname{\href{https://github.com/Deicyde/riemann-hypothesis/blob/40e56bab026552e089d7afcfeff08d2defb9c005/src/Development/AutomorphicArithmetic/AutomorphicVarieties/GaloisLibrary/KummerTheory/Mk0IdMapCongr.lean\#L197-L197}{\leanfont Kum.Xinf}}
\expandafter\newcommand\csname LN@funRep\endcsname{\href{https://github.com/Deicyde/riemann-hypothesis/blob/40e56bab026552e089d7afcfeff08d2defb9c005/src/Development/AutomorphicArithmetic/AutomorphicVarieties/GaloisLibrary/KummerTheory/Mk0IdMapCongr.lean\#L121-L124}{\leanfont Kum.funRep}}
\expandafter\newcommand\csname LN@conjCond\endcsname{\href{https://github.com/Deicyde/riemann-hypothesis/blob/40e56bab026552e089d7afcfeff08d2defb9c005/src/Development/AutomorphicArithmetic/AutomorphicVarieties/GaloisLibrary/KummerTheory/Mk0IdMapCongr.lean\#L103-L107}{\leanfont Kum.conjCond}}
\expandafter\newcommand\csname LN@EqHom\endcsname{\href{https://github.com/Deicyde/riemann-hypothesis/blob/40e56bab026552e089d7afcfeff08d2defb9c005/src/Development/AutomorphicArithmetic/AutomorphicVarieties/GaloisLibrary/KummerTheory/Mk0IdMapCongr.lean\#L109-L110}{\leanfont Kum.EqHom}}
\expandafter\newcommand\csname LN@Un\endcsname{\href{https://github.com/Deicyde/riemann-hypothesis/blob/40e56bab026552e089d7afcfeff08d2defb9c005/src/Development/AutomorphicArithmetic/AutomorphicVarieties/GaloisLibrary/KummerTheory/Mk0IdMapCongr.lean\#L186-L186}{\leanfont Kum.Un}}
\expandafter\newcommand\csname LN@funRepMem\endcsname{\href{https://github.com/Deicyde/riemann-hypothesis/blob/40e56bab026552e089d7afcfeff08d2defb9c005/src/Development/AutomorphicArithmetic/AutomorphicVarieties/GaloisLibrary/KummerTheory/FunRepMem.lean\#L5-L46}{\leanfont Kum.funRepMem}}
\expandafter\newcommand\csname LN@memPrinc\endcsname{\href{https://github.com/Deicyde/riemann-hypothesis/blob/40e56bab026552e089d7afcfeff08d2defb9c005/src/Development/AutomorphicArithmetic/AutomorphicVarieties/GaloisLibrary/KummerTheory/DinfNil.lean\#L7-L12}{\leanfont mem\_princ\_iff\_pow\_dvd\_sub\_one}}
\expandafter\newcommand\csname LN@powSucc\endcsname{\href{https://github.com/Deicyde/riemann-hypothesis/blob/40e56bab026552e089d7afcfeff08d2defb9c005/src/Development/AutomorphicArithmetic/AutomorphicVarieties/GaloisLibrary/KummerTheory/DinfNil.lean\#L14-L29}{\leanfont pow\_succ\_dvd\_pow\_prime\_sub\_one}}
\expandafter\newcommand\csname LN@qzTors\endcsname{\href{https://github.com/Deicyde/riemann-hypothesis/blob/40e56bab026552e089d7afcfeff08d2defb9c005/src/Development/AutomorphicArithmetic/AutomorphicVarieties/GaloisLibrary/KummerTheory/DinfNil.lean\#L31-L43}{\leanfont exists\_prime\_pow\_smul\_eq\_zero\_qz}}
\expandafter\newcommand\csname LN@princNhds\endcsname{\href{https://github.com/Deicyde/riemann-hypothesis/blob/40e56bab026552e089d7afcfeff08d2defb9c005/src/Development/AutomorphicArithmetic/AutomorphicVarieties/GaloisLibrary/KummerTheory/DinfNil.lean\#L45-L60}{\leanfont exists\_princ\_subset\_of\_mem\_nhds\_one}}
\expandafter\newcommand\csname LN@dinfNil\endcsname{\href{https://github.com/Deicyde/riemann-hypothesis/blob/40e56bab026552e089d7afcfeff08d2defb9c005/src/Development/AutomorphicArithmetic/AutomorphicVarieties/GaloisLibrary/KummerTheory/DinfNil.lean\#L62-L235}{\leanfont Kum.dinfNil}}
\expandafter\newcommand\csname LN@G\endcsname{\href{https://github.com/Deicyde/riemann-hypothesis/blob/40e56bab026552e089d7afcfeff08d2defb9c005/src/Formalism/ConjugateSelfDual/Statement.lean\#L6-L6}{\leanfont G}}
\expandafter\newcommand\csname LN@QZ\endcsname{\href{https://github.com/Deicyde/riemann-hypothesis/blob/40e56bab026552e089d7afcfeff08d2defb9c005/src/Formalism/GaloisCohomology/Statement.lean\#L15-L15}{\leanfont QZ}}
\expandafter\newcommand\csname LN@cycZ\endcsname{\href{https://github.com/Deicyde/riemann-hypothesis/blob/40e56bab026552e089d7afcfeff08d2defb9c005/src/ReciprocityTheory/GaloisValuations/Global.lean\#L90-L92}{\leanfont Gal.cycZ}}
\expandafter\newcommand\csname LN@Dec\endcsname{\href{https://github.com/Deicyde/riemann-hypothesis/blob/40e56bab026552e089d7afcfeff08d2defb9c005/src/ReciprocityTheory/GaloisValuations/Global.lean\#L15-L15}{\leanfont Gal.Dec}}
\expandafter\newcommand\csname LN@Ine\endcsname{\href{https://github.com/Deicyde/riemann-hypothesis/blob/40e56bab026552e089d7afcfeff08d2defb9c005/src/ReciprocityTheory/GaloisValuations/Global.lean\#L45-L64}{\leanfont Gal.Ine}}
\expandafter\newcommand\csname LN@princ\endcsname{\href{https://github.com/Deicyde/riemann-hypothesis/blob/40e56bab026552e089d7afcfeff08d2defb9c005/src/Development/AutomorphicArithmetic/CyclotomicBounds/ClassNumberBounds/TotallyRealRoots/LeNuIffPowMemPrinc.lean\#L16-L16}{\leanfont princ}}
\expandafter\newcommand\csname LN@levU\endcsname{\href{https://github.com/Deicyde/riemann-hypothesis/blob/40e56bab026552e089d7afcfeff08d2defb9c005/src/Development/AutomorphicArithmetic/AutomorphicVarieties/GaloisLibrary/SelmerTransport/ConjStableShaFam.lean\#L187-L187}{\leanfont Base.levU}}
\expandafter\newcommand\csname LN@cycU\endcsname{\href{https://github.com/Deicyde/riemann-hypothesis/blob/40e56bab026552e089d7afcfeff08d2defb9c005/src/Development/AutomorphicArithmetic/AutomorphicVarieties/GaloisLibrary/SelmerTransport/ConjStableShaFam.lean\#L189-L189}{\leanfont Base.cycU}}
\expandafter\newcommand\csname LN@conjHom\endcsname{\href{https://github.com/Deicyde/riemann-hypothesis/blob/40e56bab026552e089d7afcfeff08d2defb9c005/src/Development/AutomorphicArithmetic/AutomorphicVarieties/GaloisLibrary/SelmerTransport/RestrictionMemInvariantH1.lean\#L115-L118}{\leanfont Base.conjHom}}
\expandafter\newcommand\csname LN@conjDec\endcsname{\href{https://github.com/Deicyde/riemann-hypothesis/blob/40e56bab026552e089d7afcfeff08d2defb9c005/src/Development/AutomorphicArithmetic/AutomorphicVarieties/PartsConj.lean\#L69-L75}{\leanfont conjGal\_mem\_Dec}}
\expandafter\newcommand\csname LN@conjIne\endcsname{\href{https://github.com/Deicyde/riemann-hypothesis/blob/40e56bab026552e089d7afcfeff08d2defb9c005/src/Development/AutomorphicArithmetic/AutomorphicVarieties/PartsConj.lean\#L77-L81}{\leanfont conjGal\_mem\_Ine}}
\expandafter\newcommand\csname LN@xUnip\endcsname{\href{https://github.com/Deicyde/riemann-hypothesis/blob/40e56bab026552e089d7afcfeff08d2defb9c005/src/Development/AutomorphicArithmetic/AutomorphicVarieties/AdditiveValuationNorms/XUnip.lean\#L9-L27}{\leanfont PropA.xUnip}}
\expandafter\newcommand\csname LN@dualUnip\endcsname{\href{https://github.com/Deicyde/riemann-hypothesis/blob/40e56bab026552e089d7afcfeff08d2defb9c005/src/Development/AutomorphicArithmetic/AutomorphicVarieties/AdditiveValuationNorms/DualUnip.lean\#L9-L34}{\leanfont PropA.dualUnip}}
\expandafter\newcommand\csname LN@shaDGen\endcsname{\href{https://github.com/Deicyde/riemann-hypothesis/blob/40e56bab026552e089d7afcfeff08d2defb9c005/src/Development/AutomorphicArithmetic/AutomorphicVarieties/NumberFieldTowers/ShaDGen.lean\#L9-L59}{\leanfont Tower.shaDGen}}
\newcommand{\Pin}{40e56bab}
\newcommand{\PinUrl}{https://github.com/Deicyde/riemann-hypothesis/tree/40e56bab026552e089d7afcfeff08d2defb9c005/src}
\newcommand{\DirFiles}{64}
\newcommand{\DirLines}{5{,}460}
\newcommand{\ConcDecls}{27}
\newcommand{\ConcFiles}{13}
\newcommand{\UnitMembers}{144}
\newcommand{\UnitInDir}{144}
\newcommand{\UnitOutside}{0}
\newcommand{\UnitInNote}{13}
\newcommand{\UnitElsewhere}{131}
\newcommand{\ElsewhereRows}{The Artin map and class groups & 23\\Kummer generators, units and Selmer groups & 30\\Reflection & 8\\Restriction to the cyclotomic kernel & 32\\Definitions of the other subjects & 38\\}
\newcommand{\UnitProofs}{103}
\newcommand{\UnitSteps}{856}
\newcommand{\ConcordanceRows}{\crow{XDual}{\lean{XDual}}{Mk0IdMapCongr.lean}{90--99}
\crow{Dinf}{\lean{Dinf}}{Mk0IdMapCongr.lean}{195--195}
\crow{Xinf}{\lean{Xinf}}{Mk0IdMapCongr.lean}{197--197}
\crow{funRep}{\lean{funRep}}{Mk0IdMapCongr.lean}{121--124}
\crow{conjCond}{\lean{conjCond}}{Mk0IdMapCongr.lean}{103--107}
\crow{EqHom}{\lean{EqHom}}{Mk0IdMapCongr.lean}{109--110}
\crow{Un}{\lean{Un}}{Mk0IdMapCongr.lean}{186--186}
\crow{funRepMem}{\lean{funRepMem}}{FunRepMem.lean}{5--46}
\crow{memPrinc}{\lean{memPrinc}}{DinfNil.lean}{7--12}
\crow{powSucc}{\lean{powSucc}}{DinfNil.lean}{14--29}
\crow{qzTors}{\lean{qzTors}}{DinfNil.lean}{31--43}
\crow{princNhds}{\lean{princNhds}}{DinfNil.lean}{45--60}
\crow{dinfNil}{\lean{dinfNil}}{DinfNil.lean}{62--235}
\crow{G}{\lean{G}}{Statement.lean}{6--6}
\crow{QZ}{\lean{QZ}}{Statement.lean}{15--15}
\crow{cycZ}{\lean{cycZ}}{Global.lean}{90--92}
\crow{Dec}{\lean{Dec}}{Global.lean}{15--15}
\crow{Ine}{\lean{Ine}}{Global.lean}{45--64}
\crow{princ}{\lean{princ}}{LeNuIffPowMemPrinc.lean}{16--16}
\crow{levU}{\lean{levU}}{ConjStableShaFam.lean}{187--187}
\crow{cycU}{\lean{cycU}}{ConjStableShaFam.lean}{189--189}
\crow{conjHom}{\lean{conjHom}}{RestrictionMemInvariantH1.lean}{115--118}
\crow{conjDec}{\lean{conjDec}}{PartsConj.lean}{69--75}
\crow{conjIne}{\lean{conjIne}}{PartsConj.lean}{77--81}
\crow{xUnip}{\lean{xUnip}}{XUnip.lean}{9--27}
\crow{dualUnip}{\lean{dualUnip}}{DualUnip.lean}{9--34}
\crow{shaDGen}{\lean{shaDGen}}{ShaDGen.lean}{9--59}}
\expandafter\newcommand\csname LQ@dinfNil\endcsname{\leanquoteline{0}{theorem UpperTheory.V.Kum.dinfNil :}\leanquoteline{4}{∀ (F Ω : Type) [Field F] [NumberField F] [Field Ω] [Algebra F Ω] [IsAlgClosure F Ω] (l : ℕ) [Fact l.Prime] (N : ℕ) (U : Subgroup (G F Ω)) [U.Normal], IsClosed (U : Set (G F Ω)) → ∀ φ ∈ Dinf l N U, (∃ n, ∀ g ∈ levU l U n, funRep ℤ\_[l] (QZ l) (cycU l U) g φ = φ) ∧}\leanquoteline{4}{∀ γ ∈ levU l U 1, ∃ k : ℕ, (l : ℤ\_[l]) \textasciicircum{} k • φ = 0 ∧ ((funRep ℤ\_[l] (QZ l) (cycU l U) γ - 1) \textasciicircum{} k) φ = 0}}
\expandafter\newcommand\csname LQ@XDual\endcsname{\leanquoteline{0}{def XDual (N : ℕ) (U : Subgroup (G F Ω)) (R M : Type) [CommRing R] [AddCommGroup M] [Module R M] : Submodule R (U → M) where}\leanquoteline{2}{carrier := \{φ | IsLocallyConstant φ ∧ (∀ a b, φ (a * b) = φ a + φ b) ∧}\leanquoteline{4}{(∀ (A : ValuationSubring Ω) (u : U), u.1 ∈ Ine (F := F) A → φ u = 0) ∧}\leanquoteline{4}{∀ A : ValuationSubring Ω, A.valuation (N : Ω) < 1 → ∀ u : U, u.1 ∈ Dec (F := F) A → φ u = 0\}}\leanquoteline{2}{zero\_mem' := ⟨IsLocallyConstant.const \_, by simp, by simp, by simp⟩}\leanquoteline{2}{add\_mem' \{a b\} ha hb := ⟨ha.1.add hb.1, fun x y => by simp only [Pi.add\_apply, ha.2.1, hb.2.1]; abel,}\leanquoteline{4}{fun A u h => by simp [ha.2.2.1 A u h, hb.2.2.1 A u h],}\leanquoteline{4}{fun A hA u h => by simp [ha.2.2.2 A hA u h, hb.2.2.2 A hA u h]⟩}\leanquoteline{2}{smul\_mem' c a ha := ⟨ha.1.comp (c • ·), fun x y => by simp [ha.2.1], fun A u h => by simp [ha.2.2.1 A u h],}\leanquoteline{4}{fun A hA u h => by simp [ha.2.2.2 A hA u h]⟩}}
\expandafter\newcommand\csname LQ@funRepMem\endcsname{\leanquoteline{0}{theorem UpperTheory.V.Kum.funRepMem :}\leanquoteline{4}{∀ (F Ω : Type) [Field F] [NumberField F] [Field Ω] [Algebra F Ω] [IsAlgClosure F Ω] (N : ℕ) (U : Subgroup (G F Ω)) [U.Normal] (R M : Type) [CommRing R] [AddCommGroup M] [Module R M], (∀ g, ∀ φ ∈ XDual N U R M, funRep R M U g φ ∈ XDual N U R M) ∧ (∀ g ∈ U, ∀ φ ∈ XDual N U R M, funRep R M U g φ = φ) ∧}\leanquoteline{4}{∀ (ρ : Representation R (G F Ω) M) (φ : U → M), φ ∈ EqHom N U ρ ↔ φ ∈ XDual N U R M ∧ ∀ g u, ρ g (funRep R M U g φ u) = φ u}}
% the notations of the functions that the archive defines: one slot for each explicit argument
\newcommand{\XDual}[4]{D_{#1}(#2;#3,#4)}
\newcommand{\Dinf}[3]{D_\infty(#1,#2,#3)}
\newcommand{\Xinf}[3]{X_\infty(#1,#2,#3)}
\newcommand{\funRep}[3]{\mathcal F_{#1}(#3,#2)}
\newcommand{\conjCond}[3]{\mathcal C(#1;#2,#3)}
\newcommand{\EqHom}[3]{\mathcal E_{#1}(#2;#3)}
\newcommand{\Un}[4]{\Gal(#2/#1)_{#3,#4}}
\newcommand{\GalG}[2]{\Gal(#2/#1)}
\newcommand{\QZl}[1]{\Q_{#1}/\Z_{#1}}
\newcommand{\cyc}[4]{\kappa_{#3}^{#2/#1}(#4)}
\newcommand{\Dec}[1]{G_{#1}}
\newcommand{\Ine}[1]{I_{#1}}
\newcommand{\princ}[2]{U_{#1}^{(#2)}}
\newcommand{\levU}[3]{#2_{#1,#3}}
\newcommand{\cycU}[2]{#2_{#1,\infty}}


\linespread{1.08}                       % 11 pt on 14.7 pt
\setlength{\parindent}{0pt}
\setlength{\parskip}{5.5pt plus 1.5pt minus 0.5pt}
\setlength{\abovedisplayskip}{7.5pt plus 2pt minus 3pt}
\setlength{\belowdisplayskip}{7.5pt plus 2pt minus 3pt}
\setlength{\abovedisplayshortskip}{3pt plus 2pt}
\setlength{\belowdisplayshortskip}{6pt plus 2pt minus 2pt}
\emergencystretch=1.5em
\clubpenalty=10000 \widowpenalty=10000 \displaywidowpenalty=10000
\relpenalty=7000 \binoppenalty=9000
\raggedbottom
\pagestyle{fancy}\fancyhf{}\renewcommand{\headrulewidth}{0pt}
\fancyfoot[C]{\sffamily\fontsize{7.6}{9}\selectfont\color{inkfaint}\thepage\,/\,\pageref*{LastPage}}
\setlist{topsep=2pt,partopsep=0pt,parsep=0pt,itemsep=4pt,leftmargin=2.1em,labelsep=0.55em}

\newcommand{\Sha}{\text{\normalfont Ш}}
\newcommand{\Q}{\mathbb{Q}}
\newcommand{\Z}{\mathbb{Z}}
\newcommand{\F}{\mathbb{F}}
\newcommand{\len}{\operatorname{len}}
\newcommand{\cok}{\operatorname{cok}}
\newcommand{\End}{\operatorname{End}}
\newcommand{\Hom}{\operatorname{Hom}}
\newcommand{\Gal}{\operatorname{Gal}}
\newcommand{\diag}{\operatorname{diag}}
\newcommand{\kw}[1]{\emph{#1}}
\newcommand{\m}[1]{\mbox{$#1$}}        % a short formula that must not be broken across lines
\newcommand{\head}[1]{\par\addvspace{17pt plus 4pt minus 2pt}\Needspace{6\baselineskip}%
  \noindent{\fontsize{13.2}{16}\selectfont\bfseries\color{ink}#1}\par\nopagebreak\vspace{-1pt}}
\newcommand{\sect}[2]{\refstepcounter{section}\label{#2}\par\addvspace{19pt plus 4pt minus 2pt}\Needspace{7\baselineskip}%
  \noindent\llap{{\fontsize{13.2}{16}\selectfont\bfseries\color{inkfaint}\thesection}\hspace{9pt}}%
  {\fontsize{13.2}{16}\selectfont\bfseries\color{ink}#1}\par\nopagebreak\vspace{-1pt}}
\newcommand{\subhead}[1]{\par\addvspace{9pt plus 2pt}\Needspace{4\baselineskip}\noindent\textbf{#1}\enspace\ignorespaces}
% where a statement lives in the Lean archive
\newcommand{\leantag}{{\sffamily\bfseries\fontsize{6.3}{8}\selectfont\color{leancol}\addfontfeatures{LetterSpace=7}LEAN}}
\newcommand{\sep}{\unskip\nobreak\hspace{5pt}{\color{inkfaint}\textperiodcentered}\hspace{5pt plus 2pt}\ignorespaces}
\newcommand{\leanline}[1]{\par\nopagebreak\vspace{-2.5pt}{\fontsize{8.9}{12.4}\selectfont\color{inksoft}\raggedright
  \noindent\llap{\leantag\hspace{9pt}}#1\par}\vspace{1pt}}
% a displayed statement in a tinted box: #1 colour, #2 text
\newcommand{\statement}[2]{\par\addvspace{7pt}\begingroup\setlength{\fboxsep}{8.5pt}\setlength{\fboxrule}{0.6pt}%
  \noindent\fcolorbox{#1!30!white}{#1!3!white}{\parbox{\dimexpr\linewidth-2\fboxsep-2\fboxrule\relax}{#2}}%
  \endgroup\par\addvspace{7pt}}
\newcommand{\nt}[2]{\par\noindent\hangindent=1.1em\hangafter=1 #1\enspace{\color{inksoft}#2}\par\vspace{1pt}}
\newcommand{\legdot}[1]{\tikz[baseline=-0.62ex]\fill[#1] (0,0) circle (2.5pt);\hspace{3pt}}
\newcommand{\crow}[4]{\csname CS@#1\endcsname & #2 & \textcolor{inksoft}{#3} & \textcolor{inksoft}{#4}\\}
\newcommand{\leanquoteline}[2]{\par\noindent\hangindent=\dimexpr#1\fontdimen2\font+2.2em\relax\hangafter=1
  \hspace*{#1\fontdimen2\font}#2}
\newcommand{\quoted}[1]{\par\vspace{3pt}{\ttfamily\fontsize{7.5}{10.4}\selectfont\color{leancol}\raggedright
  \setlength{\parskip}{0pt}\leanquote{#1}\par}\vspace{4pt}}
\newcommand{\bx}[3]{{\sffamily\bfseries\fontsize{7.3}{9.2}\selectfont\color{#1!82!black}#2}\\[1.5pt]#3}
\newcommand{\cs}[2]{\expandafter\newcommand\csname CS@#1\endcsname{#2}}

% theorems
\newtheoremstyle{plainx}{7pt plus 2pt}{7pt plus 2pt}{\itshape}{}{\bfseries}{.}{0.5em}{}
\newtheoremstyle{uprightx}{7pt plus 2pt}{7pt plus 2pt}{}{}{\bfseries}{.}{0.5em}{}
\theoremstyle{plainx}
\newtheorem{theorem}{Theorem}[section]
\newtheorem{lemma}[theorem]{Lemma}
\newtheorem{proposition}[theorem]{Proposition}
\newtheorem{corollary}[theorem]{Corollary}
\theoremstyle{uprightx}
\newtheorem{definition}[theorem]{Definition}
\newtheorem{setting}[theorem]{Setting}
\newtheorem{remark}[theorem]{Remark}
\newtheorem{example}[theorem]{Example}
% the ambient objects of the definitions, said once (implicit and instance arguments of the archive; no hypothesis)
\newenvironment{conventions}{\par\addvspace{5pt}\noindent\textbf{Conventions.}\enspace\ignorespaces}{\par\addvspace{3pt}}
\renewcommand{\qedsymbol}{{\setlength{\fboxsep}{0pt}\setlength{\fboxrule}{0.5pt}\fbox{\rule{0pt}{0.46em}\hspace{0.46em}}}}
\makeatletter
\renewenvironment{proof}[1][Proof]{\par\pushQED{\qed}\normalfont\topsep0pt\partopsep0pt
  \trivlist\item[\hskip\labelsep\itshape #1.]\ignorespaces}{\popQED\endtrivlist\@endpefalse}
\makeatother

% where each declaration of the concordance is treated
\cs{XDual}{Def.~\ref{d:char}} \cs{Dinf}{Def.~\ref{d:char}} \cs{Xinf}{Def.~\ref{d:char}}
\cs{funRep}{Def.~\ref{d:conj}} \cs{conjCond}{Def.~\ref{d:equiv}} \cs{EqHom}{Def.~\ref{d:equiv}}
\cs{Un}{Def.~\ref{d:groups}} \cs{funRepMem}{Prop.~\ref{p:action}} \cs{memPrinc}{Lemma~\ref{l:four}(a)}
\cs{powSucc}{Lemma~\ref{l:four}(b)} \cs{qzTors}{Lemma~\ref{l:four}(c)} \cs{princNhds}{Lemma~\ref{l:four}(d)}
\cs{dinfNil}{Theorem A} \cs{G}{Def.~\ref{d:galois}} \cs{QZ}{Def.~\ref{d:galois}} \cs{cycZ}{Def.~\ref{d:galois}}
\cs{Dec}{Def.~\ref{d:char}} \cs{Ine}{Def.~\ref{d:char}} \cs{princ}{Def.~\ref{d:groups}}
\cs{levU}{Def.~\ref{d:groups}} \cs{cycU}{Def.~\ref{d:groups}} \cs{conjHom}{Def.~\ref{d:conj}}
\cs{conjDec}{Prop.~\ref{p:action}, proof} \cs{conjIne}{Prop.~\ref{p:action}, proof} \cs{xUnip}{Cor.~\ref{c:unip}}
\cs{dualUnip}{Prop.~\ref{p:dual}} \cs{shaDGen}{Prop.~\ref{p:inv}}

\begin{document}
\color{ink}

% ------------------------------------------------------------------------------------- title
{\raggedright
{\sffamily\bfseries\fontsize{7.2}{9}\selectfont\color{leancol}\addfontfeatures{LetterSpace=8}DEFORMALIZATION 2\par}
\vspace{5pt}
{\fontsize{21.5}{25.5}\selectfont\bfseries Local nilpotence over a cyclotomic tower\par}
\vspace{5pt}
{\fontsize{12.4}{16}\selectfont\itshape\color{inksoft} One part of a unit of the Lean archive behind the proof of the
Riemann hypothesis announced by L.\,Alpöge, read back into mathematics\par}}
\vspace{8pt}
{\color{rulegrey}\hrule height 0.6pt}
\vspace{7pt}
\thispagestyle{empty}

\textbf{Summary.} Let $F$ be a number field with algebraic closure~$\Omega$, let $\ell$ be a prime, and let
$\cyc{F}{\Omega}{\ell}{\cdot}$ be the $\ell$-adic cyclotomic character of the Galois group~$\GalG{F}{\Omega}$. This
group acts by conjugation on the characters of the kernel of the cyclotomic character, with values
in~$\QZl{\ell}$. We show that every locally constant character~$\varphi$ is fixed by all~$g$ with
\m{\cyc{F}{\Omega}{\ell}{g}\equiv1} modulo~$\ell^n$, for some~$n$, and that for every $\gamma$ with
\m{\cyc{F}{\Omega}{\ell}{\gamma}\equiv1} modulo~$\ell$ a power of the operator ``action of~$\gamma$, minus the
identity'' kills~$\varphi$; the exponent $\ell^nk_0$ suffices, where $\ell^{k_0}$ kills~$\varphi$. This is the
\kw{local nilpotence} that the first note of this series~\cite{N1} used as an input. By duality it makes the
action of~$\gamma$ unipotent modulo~$\ell$ on the Iwasawa module. The proof is a compactness argument in the
profinite group~$\GalG{F}{\Omega}$ and one congruence between polynomials.

\statement{leancol}{\fontsize{9.6}{13.2}\selectfont\textbf{What this note is.} On 20 November 2026 L.\,Alpöge
announced a proof of the Riemann hypothesis, with a Lean archive \cite{Al}. That is a claim, and nothing here asserts
that the hypothesis is proved. The reading project divides the part of the archive that the final theorems use
into units. This note treats one part of one unit, on Kummer theory: \UnitInNote{} of its \UnitMembers{}
declarations, those that define the modules of characters and prove local nilpotence. The other \UnitElsewhere{}
belong to notes to come (§\ref{s:notes}). Everything stated in §§\ref{s:intro}--\ref{s:proof} is proved in full;
§\ref{s:use} states the two results of the archive that use the theorem, one derived here from a proposition
of the first note, the other quoted without proof. Every statement was compared
with its formal counterpart at commit \href{\PinUrl}{{\leanfont\Pin}} of the public mirror; Appendix~A is the
concordance of \ConcDecls{} declarations. The note was written from the structured proofs of the reading project,
without opening the Lean files of the unit (§\ref{s:notes}).}

% ------------------------------------------------------------------------------------- 1
\sect{The statement}{s:intro}
In the first note~\cite{N1} the lengths of certain groups of Galois cohomology along the cyclotomic tower of a
number field~$F$ were expressed through one operator on one module, the Iwasawa module
\m{\Xinf{\ell}{\ell}{\GalG{F}{\Omega}}} of a prime~$\ell$ and an algebraic closure~$\Omega$ of~$F$, and the
estimates needed that this operator is unipotent modulo~$\ell$. This module is the dual of a discrete module
\m{\Dinf{\ell}{\ell}{\GalG{F}{\Omega}}} of characters, and on a discrete module unipotence takes the form proved
here: every element is killed by a power of the operator minus the identity. The first note stated this as its
Input~1, with a sketch. This note proves it, in the generality of the archive: there the second argument of the
two modules is any integer \m{N\ge0}, and the third is any closed normal subgroup of~$\GalG{F}{\Omega}$. The
notation is set up in §\ref{s:setting}.

\statement{arith}{\phantomsection\label{t:A}\textbf{Theorem A} (local nilpotence).\enspace\emph{Let $\ell$ be a prime
number, $F$ a number field and $\Omega$ an algebraic closure of~$F$. Let $H$ be a closed normal subgroup
of~$\GalG{F}{\Omega}$, let \m{N\ge0} be an integer, and let $\varphi$ be an element of the module
$\Dinf{\ell}{N}{H}$ of characters of~$\cycU{\ell}{H}$. Write}
\[
\pi=\funRep{\Z_\ell}{\QZl{\ell}}{\cycU{\ell}{H}}
\]
\emph{for the representation of $\GalG{F}{\Omega}$ by conjugation on the functions from $\cycU{\ell}{H}$
to~$\QZl{\ell}$.}
\begin{enumerate}[label=\textup{(\roman*)},topsep=3pt,itemsep=1pt]
\item \emph{There is an integer \m{n\ge0} such that \m{\pi(g)\varphi=\varphi} for every
\m{g\in\levU{\ell}{H}{n}}.}
\item \emph{Some power of $\ell$ kills~$\varphi$. If \m{\ell^{k_0}\varphi=0} and $n$ is as in \textup{(i)}, then for
every \m{\gamma\in\levU{\ell}{H}{1}} and \m{k=\ell^nk_0}}
\[
(\pi(\gamma)-1)^{k}\varphi=0\qquad\text{\emph{and}}\qquad\ell^{k}\varphi=0 .
\]
\end{enumerate}}
\leanline{\lean{dinfNil}}

Here $\GalG{F}{\Omega}$ is the Galois group of $\Omega$ over~$F$ with its usual topology,
\m{\cyc{F}{\Omega}{\ell}{g}} is the value at~$g$ of its $\ell$-adic cyclotomic character, $\levU{\ell}{H}{n}$ is the
group of all \m{g\in H} with \m{\cyc{F}{\Omega}{\ell}{g}\equiv1} modulo~$\ell^n$, and $\cycU{\ell}{H}$ is the
group of all \m{g\in H} with \m{\cyc{F}{\Omega}{\ell}{g}=1}. The group $\QZl{\ell}$ is the quotient of the
$\ell$-adic numbers by the $\ell$-adic integers, the module $\Dinf{\ell}{N}{H}$ is that of
Definition~\ref{d:char}, and $\pi(g)\varphi$ is the character \m{x\mapsto\varphi(g^{-1}xg)}. The archive
states (ii) in this form: for every~$\gamma$ there is an exponent~$k$ with the two properties. Its proof gives the
value above, the same for all~$\gamma$.

% ------------------------------------------------------------------------------------- 2
\sect{Characters and conjugation}{s:setting}
\begin{conventions}\phantomsection\label{c:conventions}%
In the definitions of this note $\ell$ is a prime number, $\Z_\ell$ is the ring of $\ell$-adic integers and $\Q_\ell$
its field of fractions; $F$ is a number field, and $\Omega$ is an algebraic closure of~$F$. A notation shows the
prime and the two fields exactly where the formal definition has them as written arguments; elsewhere they are
understood. A theorem names them again in its own statement.
\end{conventions}

\begin{definition}[the Galois group, the coefficients, the cyclotomic character]\label{d:galois}
The \kw{Galois group} $\GalG{F}{\Omega}$ is the group of the automorphisms of the field~$\Omega$ that fix every
element of~$F$. It carries its usual topology: the subgroups~$\GalG{E}{\Omega}$, for $E$ a finite extension of $F$
in~$\Omega$, are a basis of neighbourhoods of~$1$.

The group $\QZl{\ell}$ is the quotient of the additive group of~$\Q_\ell$ by~$\Z_\ell$. It is a module over~$\Z_\ell$.

The \kw{$\ell$-adic cyclotomic character} of $\Omega$ over~$F$ is the homomorphism
\m{g\mapsto\cyc{F}{\Omega}{\ell}{g}} from $\GalG{F}{\Omega}$ to the group $\Z_\ell^\times$ of units of~$\Z_\ell$ such
that \m{g(\zeta)=\zeta^{\cyc{F}{\Omega}{\ell}{g}}} for every \m{g\in\GalG{F}{\Omega}} and every root of unity $\zeta$
of $\ell$-power order in~$\Omega$. It is continuous.
\end{definition}
\leanline{\lean{G}\sep\lean{QZ}\sep\lean{cycZ}}

The group $\GalG{F}{\Omega}$ is profinite: it is compact, and every neighbourhood of~$1$ contains an open normal
subgroup \cite[Ch.\,I, §1]{NSW}.

\begin{definition}[the groups]\label{d:groups}
For an integer \m{n\ge1} let \m{\princ{\ell}{n}=1+\ell^n\Z_\ell}, and let \m{\princ{\ell}{0}=\Z_\ell^\times}; these are
open subgroups of~$\Z_\ell^\times$. For a subgroup $H$ of~$\GalG{F}{\Omega}$ and an integer \m{n\ge0} put
\[
\levU{\ell}{H}{n}=\bigl\{g\in H:\ \cyc{F}{\Omega}{\ell}{g}\in\princ{\ell}{n}\bigr\},\qquad
\cycU{\ell}{H}=\bigl\{g\in H:\ \cyc{F}{\Omega}{\ell}{g}=1\bigr\}.
\]
For the whole group we also write \m{\Un{F}{\Omega}{\ell}{n}} for the group of all \m{g\in\GalG{F}{\Omega}} with
\m{\cyc{F}{\Omega}{\ell}{g}\in\princ{\ell}{n}}; it is the first of the two groups for \m{H=\GalG{F}{\Omega}}.
\end{definition}
\leanline{\lean{princ}\sep\lean{levU}\sep\lean{cycU}\sep\lean{Un}}

So \m{\levU{\ell}{H}{0}=H}, the groups $\levU{\ell}{H}{n}$ decrease as $n$ grows, and $\cycU{\ell}{H}$, the
intersection of $H$ with the kernel of the cyclotomic character, is their intersection. If $H$ is normal
in~$\GalG{F}{\Omega}$, so are all of these. The group $\Un{F}{\Omega}{\ell}{n}$ is the Galois group of $\Omega$ over
the field generated over~$F$ by the roots of unity of order~$\ell^n$, and $\cycU{\ell}{\GalG{F}{\Omega}}$ is that of
$\Omega$ over the field $F_\infty$ generated by all roots of unity of $\ell$-power order.

\begin{definition}[characters]\label{d:char}
A \kw{valuation ring} of $\Omega$ is a subring $A$ such that \m{z\in A} or \m{z^{-1}\in A} for every non-zero
\m{z\in\Omega}; its maximal ideal $\mathfrak m_A$ consists of the non-units of~$A$. The \kw{decomposition group}
$\Dec{A}$ is the group of all \m{g\in\GalG{F}{\Omega}} with \m{g(A)=A}, and the \kw{inertia group} $\Ine{A}$ is the
subgroup of those that act trivially on the residue field~\m{A/\mathfrak m_A}.

Let $C$ be a subgroup of~$\GalG{F}{\Omega}$, let \m{N\ge0} be an integer, and let $M$ be a module over a commutative
ring~$R$. The module \m{\XDual{N}{C}{R}{M}} consists of the homomorphisms \m{\varphi\colon C\to M} that are locally
constant, vanish on \m{C\cap\Ine{A}} for every valuation ring $A$ of~$\Omega$, and vanish on \m{C\cap\Dec{A}} for
every valuation ring $A$ with \m{N\in\mathfrak m_A}. It is an $R$-module under pointwise operations. For a subgroup
$H$ of~$\GalG{F}{\Omega}$ put
\[
\Dinf{\ell}{N}{H}=\XDual{N}{\cycU{\ell}{H}}{\Z_\ell}{\QZl{\ell}},\qquad
\Xinf{\ell}{N}{H}=\Hom_{\Z_\ell}\bigl(\Dinf{\ell}{N}{H},\QZl{\ell}\bigr).
\]
\end{definition}
\leanline{\lean{Dec}\sep\lean{Ine}\sep\lean{XDual}\sep\lean{Dinf}\sep\lean{Xinf}}

The elements of \m{\XDual{N}{C}{R}{M}} are the characters of~$C$ that are unramified at every place of~$\Omega$ and
trivial on the decomposition groups of the places above the primes that divide~$N$. For \m{H=\GalG{F}{\Omega}} and
\m{N=\ell} the modules \m{\Dinf{\ell}{\ell}{\GalG{F}{\Omega}}} and \m{\Xinf{\ell}{\ell}{\GalG{F}{\Omega}}} are those of
the first note \cite[§7]{N1}, with \m{L=F}.

\begin{definition}[conjugation]\label{d:conj}
Let $C$ be a normal subgroup of~$\GalG{F}{\Omega}$ and \m{g\in\GalG{F}{\Omega}}. \kw{Conjugation} by~$g$ is the
homomorphism from $C$ to~$C$ that sends \m{x\in C} to~\m{g^{-1}xg}. Let $M$ be a module over a commutative ring~$R$.
We write \m{\funRep{R}{M}{C}} for the following representation of the group $\GalG{F}{\Omega}$ on the $R$-module of
all functions \m{C\to M}: the element~$g$ sends a function~$\varphi$ to the function
\m{\funRep{R}{M}{C}(g)\varphi} with
\[
\bigl(\funRep{R}{M}{C}(g)\varphi\bigr)(x)=\varphi(g^{-1}xg)\qquad(x\in C).
\]
Each \m{\funRep{R}{M}{C}(g)} is $R$-linear, and for all \m{g,h\in\GalG{F}{\Omega}}
\[
\funRep{R}{M}{C}(gh)=\funRep{R}{M}{C}(g)\,\funRep{R}{M}{C}(h).
\]
\end{definition}
\leanline{\lean{conjHom}\sep\lean{funRep}}

The archive defines conjugation for a normal subgroup of any group; only the group $\GalG{F}{\Omega}$ occurs in
this note.

\begin{definition}[equivariant characters]\label{d:equiv}
Let $C$ be a subgroup of~$\GalG{F}{\Omega}$ and let $M$ be a module over a commutative ring~$R$. For an
automorphism~$\sigma$ of the field~$\Omega$ and an $R$-linear map \m{\theta\colon M\to M}, the module
\m{\conjCond{C}{\sigma}{\theta}} consists of the functions \m{\varphi\colon C\to M} such that
\[
\varphi(x')=\theta\bigl(\varphi(x)\bigr)\qquad\text{for all }x,x'\in C\text{ with }x'\circ\sigma=\sigma\circ x .
\]
It is an $R$-module under pointwise operations. For an integer \m{N\ge0} and a representation~$\rho$ of
$\GalG{F}{\Omega}$ on~$M$ by $R$-linear maps put
\[
\EqHom{N}{C}{\rho}=\XDual{N}{C}{R}{M}\ \cap\bigcap_{g\in\GalG{F}{\Omega}}\conjCond{C}{g}{\rho(g)} .
\]
\end{definition}
\leanline{\lean{conjCond}\sep\lean{EqHom}}

For an element $g$ of~$\GalG{F}{\Omega}$ in the place of~$\sigma$, the relation \m{x'\circ g=g\circ x} says that
\m{x'=gxg^{-1}}. So a function~$\varphi$ lies in \m{\EqHom{N}{C}{\rho}} if and only if it lies in
\m{\XDual{N}{C}{R}{M}} and
\[
\varphi(gxg^{-1})=\rho(g)\varphi(x)\qquad\text{for all }g\in\GalG{F}{\Omega}\text{ and all }x\in C\text{ with }gxg^{-1}\in C .
\]

\begin{proposition}\label{p:action}
Let $F$ be a number field and $\Omega$ an algebraic closure of~$F$. Let $C$ be a normal subgroup
of~$\GalG{F}{\Omega}$, let \m{N\ge0} be an integer, and let $M$ be a module over a commutative ring~$R$.
\begin{enumerate}[label=\textup{(\alph*)},topsep=2pt,itemsep=2pt]
\item If \m{\varphi\in\XDual{N}{C}{R}{M}} and \m{g\in\GalG{F}{\Omega}}, then
\m{\funRep{R}{M}{C}(g)\varphi\in\XDual{N}{C}{R}{M}}.
\item If \m{\varphi\in\XDual{N}{C}{R}{M}} and \m{g\in C}, then \m{\funRep{R}{M}{C}(g)\varphi=\varphi}.
\item Let $\rho$ be a representation of $\GalG{F}{\Omega}$ on~$M$ by $R$-linear maps. A function
\m{\varphi\colon C\to M} lies in \m{\EqHom{N}{C}{\rho}} if and only if it lies in \m{\XDual{N}{C}{R}{M}} and
\m{\rho(g)\bigl((\funRep{R}{M}{C}(g)\varphi)(x)\bigr)=\varphi(x)} for all \m{g\in\GalG{F}{\Omega}} and \m{x\in C}.
\end{enumerate}
\end{proposition}
\begin{proof}
(a) The map \m{x\mapsto g^{-1}xg} is a continuous homomorphism of $C$ to itself, so
\m{\funRep{R}{M}{C}(g)\varphi} is a locally constant homomorphism. Let $A$ be a valuation ring of~$\Omega$. Then
\m{g^{-1}(A)} is a valuation ring with maximal ideal \m{g^{-1}(\mathfrak m_A)}, and
\[
g^{-1}\Dec{A}g=\Dec{g^{-1}(A)},\qquad g^{-1}\Ine{A}g=\Ine{g^{-1}(A)} .
\]
So for \m{x\in C\cap\Ine{A}} the element \m{g^{-1}xg} lies in \m{C\cap\Ine{g^{-1}(A)}}, where $\varphi$ vanishes. If
\m{N\in\mathfrak m_A}, then \m{N=g^{-1}(N)} lies in \m{g^{-1}(\mathfrak m_A)}, and the same argument applies to the
decomposition groups.

(b) For \m{g\in C} the three factors of \m{g^{-1}xg} lie in~$C$, and
\m{\varphi(g^{-1}xg)=-\varphi(g)+\varphi(x)+\varphi(g)=\varphi(x)}, because $M$ is commutative.

(c) As $C$ is normal, the condition displayed after Definition~\ref{d:equiv} is a condition for all
\m{g\in\GalG{F}{\Omega}} and all \m{x\in C}. Put \m{y=g^{-1}xg}. Then the condition of~(c) reads
$\rho(g)\varphi(y)=\varphi(gyg^{-1})$, and $y$ runs through $C$ when $x$ does.
\end{proof}
\leanline{\lean{funRepMem}, with \lean{conjDec} and \lean{conjIne} in the proof of (a)\sep the module of (c):
\lean{EqHom}, defined through \lean{conjCond}}

By (a) the group $\GalG{F}{\Omega}$ acts on \m{\XDual{N}{C}{R}{M}}, and by (b) the action factors through the
quotient of $\GalG{F}{\Omega}$ by~$C$. For \m{C=\cycU{\ell}{H}}, \m{R=\Z_\ell} and \m{M=\QZl{\ell}} this is the action
on $\Dinf{\ell}{N}{H}$ in Theorem~A, and it factors through the quotient of $\GalG{F}{\Omega}$ by~$\cycU{\ell}{H}$.

% ------------------------------------------------------------------------------------- 3
\sect{Four facts about \texorpdfstring{$\ell$}{l}-adic numbers}{s:four}
\begin{lemma}\label{l:four}
Let $\ell$ be a prime number.
\begin{enumerate}[label=\textup{(\alph*)},topsep=2pt,itemsep=2pt]
\item Let \m{m\ge0} be an integer. A unit $u$ of~$\Z_\ell$ lies in the kernel of the reduction of units
modulo~$\ell^m$ if and only if $\ell^m$ divides \m{u-1}. So this kernel is~$\princ{\ell}{m}$.
\item Let \m{v\in\Z_\ell} and \m{j\ge0}. If $\ell^{j+1}$ divides \m{v-1}, then $\ell^{j+2}$ divides \m{v^\ell-1}.
\item Every element of~$\QZl{\ell}$ is killed by a power of~$\ell$.
\item Every neighbourhood of~$1$ in~$\Z_\ell^\times$ contains $\princ{\ell}{n}$ for some~$n$.
\end{enumerate}
\end{lemma}
\begin{proof}
(a) The reduction of $u$ is~$1$ exactly when \m{u-1} lies in the kernel $\ell^m\Z_\ell$ of the reduction of~$\Z_\ell$.

(b) Write \m{v-1=\ell^{j+1}r} with \m{r\in\Z_\ell}, and \m{v^\ell-1=(v-1)\,s} with \m{s=1+v+\dots+v^{\ell-1}}. Each
\m{v^i-1} is divisible by \m{v-1}, so \m{s-\ell=\sum_{i<\ell}(v^i-1)=(v-1)\,d} for some \m{d\in\Z_\ell}. Hence
\[
v^\ell-1=\ell^{j+1}r\,\bigl(\ell+\ell^{j+1}r\,d\bigr)=\ell^{j+2}\bigl(r+\ell^{j}r^2d\bigr).
\]

(c) An element of $\QZl{\ell}$ is the class of some \m{r\in\Q_\ell}, and \m{\ell^mr\in\Z_\ell} for large~$m$.

(d) The topology of $\Z_\ell^\times$ is induced from~$\Z_\ell$, so a neighbourhood of~$1$ contains all units $u$
with \m{u-1\in\ell^n\Z_\ell} for some~$n$.
\end{proof}
\leanline{(a) \lean{memPrinc}\sep (b) \lean{powSucc}\sep (c) \lean{qzTors}\sep (d) \lean{princNhds}}

By (b) and induction, an element $v$ of \m{1+\ell\Z_\ell} satisfies
\begin{equation}\label{e:pow}
v^{\ell^j}\equiv1\pmod{\ell^{j+1}}\qquad\text{for every }j\ge0 .
\end{equation}

% ------------------------------------------------------------------------------------- 4
\sect{Proof of Theorem A}{s:proof}
We are in the situation of Theorem~A. Write \m{C=\cycU{\ell}{H}}; then the representation of the theorem is
\m{\pi=\funRep{\Z_\ell}{\QZl{\ell}}{C}}. The group $C$ is a normal subgroup of~$\GalG{F}{\Omega}$, and it is closed,
because $H$ is closed and the cyclotomic character \m{g\mapsto\cyc{F}{\Omega}{\ell}{g}} is continuous. So $C$ is
compact. The proof uses of $\varphi$ only that it is a locally constant homomorphism \m{C\to\QZl{\ell}}.

\subhead{Proof of (i).} The kernel of~$\varphi$ is an open subgroup of~$C$, so it is the intersection of $C$ with
an open subset of~$\GalG{F}{\Omega}$ that contains~$1$. This subset contains an open normal subgroup $W$
of~$\GalG{F}{\Omega}$. Thus
\begin{equation}\label{e:w}
\varphi(x)=0\qquad\text{for all }x\in W\cap C .
\end{equation}
\emph{Every \m{c\in C} fixes~$\varphi$}, that is, \m{\pi(c)\varphi=\varphi}. This is Proposition~\ref{p:action}(b);
it uses only that $\varphi$ is a homomorphism.

\emph{Every \m{w\in W} fixes~$\varphi$.} Let \m{x\in C}. The element \m{w^{-1}xwx^{-1}} lies in~$C$, because $C$ is
normal, and in~$W$, because it is the product of $w^{-1}$ and \m{xwx^{-1}} and $W$ is normal. By \eqref{e:w},
\[
(\pi(w)\varphi)(x)=\varphi\bigl(w^{-1}xwx^{-1}\cdot x\bigr)=\varphi\bigl(w^{-1}xwx^{-1}\bigr)+\varphi(x)=\varphi(x).
\]
Since \m{\pi(ab)\varphi=\pi(a)(\pi(b)\varphi)} for \m{a,b\in\GalG{F}{\Omega}}, every element $g$ of the set
\[
S=\{wc:\ w\in W,\ c\in C\}
\]
fixes~$\varphi$: \m{\pi(g)\varphi=\varphi}. The set $S$ is a union of cosets of the open subgroup~$W$, so it is open.

Let $K$ be the set of the elements of $H$ that are not in~$S$. It is closed in the compact
group~$\GalG{F}{\Omega}$, hence compact, and the set \m{\cyc{F}{\Omega}{\ell}{K}} of the values
\m{\cyc{F}{\Omega}{\ell}{g}}, \m{g\in K}, is a compact, hence closed, subset of~$\Z_\ell^\times$. It does not
contain~$1$: an element $g$ of $H$ with \m{\cyc{F}{\Omega}{\ell}{g}=1} lies in $C$, and \m{C\subseteq S}. By
Lemma~\ref{l:four}(d) there is $n$ such that $\princ{\ell}{n}$ does not meet~\m{\cyc{F}{\Omega}{\ell}{K}}. Let
\m{g\in\levU{\ell}{H}{n}}. Then \m{\cyc{F}{\Omega}{\ell}{g}\in\princ{\ell}{n}}, so \m{g\notin K}; as \m{g\in H}, this
means \m{g\in S}, and \m{\pi(g)\varphi=\varphi}.

\subhead{Proof of (ii).} The function $\varphi$ is locally constant on a compact space, so it takes finitely many
values. By Lemma~\ref{l:four}(c) each of them is killed by a power of~$\ell$. So \m{\ell^{k_0}\varphi=0} for some
\m{k_0\ge0}.

Let $n$ be as in~(i), put \m{q=\ell^n}, and let \m{\gamma\in\levU{\ell}{H}{1}}. Then
\m{\cyc{F}{\Omega}{\ell}{\gamma}\in1+\ell\Z_\ell}, and by \eqref{e:pow} with \m{j=n} the unit
\m{\cyc{F}{\Omega}{\ell}{\gamma^q}=\cyc{F}{\Omega}{\ell}{\gamma}^{q}} lies in~$\princ{\ell}{n}$. So
\m{\gamma^q\in\levU{\ell}{H}{n}}, and by (i), since \m{\pi(\gamma)^q=\pi(\gamma^q)},
\begin{equation}\label{e:q}
\pi(\gamma)^q\varphi=\varphi .
\end{equation}
Let $T$ be the endomorphism \m{\pi(\gamma)-1} of the $\Z_\ell$-module of all functions \m{C\to\QZl{\ell}}, and let
$J$ be the set of polynomials \m{P\in\Z[Y]} with \m{P(T)\varphi=0}. It is an ideal of the polynomial ring~$\Z[Y]$.
It contains
\[
(Y+1)^q-1\qquad\text{and}\qquad\ell^{k_0},
\]
the first by \eqref{e:q}, since \m{T+1=\pi(\gamma)}, and the second by the choice of~$k_0$. Now $q$ is a power
of~$\ell$, so \m{(Y+1)^q=Y^q+1} in the polynomial ring over the field with $\ell$ elements. Hence there is
\m{Q\in\Z[Y]} with
\[
(Y+1)^q-Y^q-1=\ell\,Q ,\qquad\text{that is,}\qquad Y^q=\bigl((Y+1)^q-1\bigr)-\ell\,Q .
\]
Modulo~$J$ this gives \m{Y^q\equiv-\ell Q}, and then
\[
Y^{qk_0}\equiv(-Q)^{k_0}\,\ell^{k_0}\equiv0\pmod J .
\]
So \m{T^{qk_0}\varphi=0}, which is \m{(\pi(\gamma)-1)^k\varphi=0} for \m{k=qk_0=\ell^nk_0}. Finally
\m{\ell^k\varphi=0}, as \m{k\ge k_0}. \hfill\qedsymbol

\medskip
The exponent does not depend on~$\gamma$. For \m{\ell=2} every unit of $\Z_2$ is congruent to~$1$ modulo~$2$, so
\m{\levU{2}{H}{1}=H} and (ii) holds for every \m{\gamma\in H}.

% ------------------------------------------------------------------------------------- 5
\sect{Where the theorem is used}{s:use}
The archive uses Theorem~A in two places outside the unit of this note. Both are stated here as the archive
states them. The first is derived from Theorem~A(ii) and a proposition that the first note proves; the second
uses Theorem~A(i) and is not proved here.

\subhead{Unipotence on the dual.} An endomorphism~$t$ of a module over~$\Z_\ell$ is \kw{unipotent modulo~$\ell$}
if some power of \m{t-1} maps the module into its multiples by~$\ell$.

\begin{proposition}[the dual module]\label{p:dual}
Let $\ell$ be a prime number, let $B$ be a module over~$\Z_\ell$ and $f$ an endomorphism of~$B$. Let $t$ be an
endomorphism of the $\Z_\ell$-module \m{\Hom_{\Z_\ell}(B,\QZl{\ell})} such that
\m{(t\xi)(\beta)=\xi(f\beta)} for all \m{\xi\in\Hom_{\Z_\ell}(B,\QZl{\ell})} and all \m{\beta\in B}. Suppose that
the module \m{\Hom_{\Z_\ell}(B,\QZl{\ell})} is finitely generated over~$\Z_\ell$, and that every element of~$B$ is
killed by a power of \m{f-1}. Then $t$ is unipotent modulo~$\ell$.
\end{proposition}
\leanline{\lean{dualUnip}}

This is the proposition on the dual module of the first note \cite[§7]{N1}, where it is proved.

\begin{corollary}[unipotence on the dual]\label{c:unip}
Let $\ell$ be a prime number, $F$ a number field and $\Omega$ an algebraic closure of~$F$. Let $H$ be a closed
normal subgroup of~$\GalG{F}{\Omega}$ and \m{N\ge0} an integer such that the module $\Xinf{\ell}{N}{H}$ is
finitely generated over~$\Z_\ell$, and write \m{\pi=\funRep{\Z_\ell}{\QZl{\ell}}{\cycU{\ell}{H}}}. Let
\m{\gamma\in\levU{\ell}{H}{1}}, and let $t$ be the action of~$\gamma$ on~$\Xinf{\ell}{N}{H}$:
\[
(t\xi)(\varphi)=\xi\bigl(\pi(\gamma^{-1})\varphi\bigr)\qquad\text{for }\xi\in\Xinf{\ell}{N}{H},\ \varphi\in\Dinf{\ell}{N}{H}.
\]
Then $t$ is unipotent modulo~$\ell$.
\end{corollary}
\begin{proof}
Put \m{B=\Dinf{\ell}{N}{H}}. By Proposition~\ref{p:action}(a) the operator \m{\pi(\gamma^{-1})} maps $B$ to itself;
let $f$ be its restriction to~$B$. The element $\gamma^{-1}$ lies in~$\levU{\ell}{H}{1}$. By Theorem~A(ii)
for~$\gamma^{-1}$, the operator \m{f-1} is locally nilpotent on~$B$: every element is killed by a power of it. The
module \m{\Hom_{\Z_\ell}(B,\QZl{\ell})} is~$\Xinf{\ell}{N}{H}$, which is finitely generated over~$\Z_\ell$, and
\m{(t\xi)(\beta)=\xi(f\beta)}. Proposition~\ref{p:dual} applies.
\end{proof}
\leanline{\lean{xUnip}, by \lean{dualUnip}}

For \m{H=\GalG{F}{\Omega}} and \m{N=\ell} this is how Input~1 of the first note enters the proof of its Theorem~D.

\subhead{Invariants.} Part (i) of Theorem~A is used in the descent to a layer of the tower
\cite[§8, Input~2b]{N1}, through the following statement.

\begin{proposition}[invariants]\label{p:inv}
Let $\ell$ be a prime number, $F$ a number field and $\Omega$ an algebraic closure of~$F$. Let $H$ be a closed
normal subgroup of~$\GalG{F}{\Omega}$, let \m{N\ge0} be an integer and $e$ an integer, and write
\m{\pi=\funRep{\Z_\ell}{\QZl{\ell}}{\cycU{\ell}{H}}}. Let $V$ be a subgroup of~$\GalG{F}{\Omega}$ contained in~$H$,
and let \m{\eta\in V} be such that every value \m{\cyc{F}{\Omega}{\ell}{g}}, \m{g\in V}, lies in the closure
in~$\Z_\ell^\times$ of the set of the powers \m{\cyc{F}{\Omega}{\ell}{\eta}^{i}}, \m{i\in\Z}. Then for an element
$\varphi$ of~$\Dinf{\ell}{N}{H}$ the following three conditions are equivalent:
\begin{enumerate}[label=\textup{(\arabic*)},topsep=2pt,itemsep=2pt]
\item \m{\cyc{F}{\Omega}{\ell}{g}^{-e}\,\pi(g)\varphi=\varphi} for every \m{g\in V};
\item \m{\cyc{F}{\Omega}{\ell}{\eta}^{-e}\,\pi(\eta)\varphi=\varphi};
\item \m{\cyc{F}{\Omega}{\ell}{\eta^{-1}}^{-e}\,\pi(\eta^{-1})\varphi=\varphi}.
\end{enumerate}
\end{proposition}
\leanline{\lean{shaDGen}}

The archive states this as two equalities of submodules of~$\Dinf{\ell}{N}{H}$: the elements with~(1) are the
kernel of the operator \m{\cyc{F}{\Omega}{\ell}{\eta}^{-e}\,\pi(\eta)-1}, and they are the kernel of the same
operator for~$\eta^{-1}$. For \m{H=\GalG{F}{\Omega}} it is the statement behind Input~2b of the first note. It
belongs to another unit of the archive and is not proved in this note; we only say where Theorem~A enters. One
needs that an element $\varphi$ of~$\Dinf{\ell}{N}{H}$ with~(2) has~(1): that an element fixed by~$\eta$ is fixed
by~$V$, for the action twisted by a power of the cyclotomic character. By Theorem~A(i) the element is fixed
by~$\levU{\ell}{H}{n}$ for some~$n$. This also holds for the twisted action when $n$ is at least the
exponent~$k_0$ of Theorem~A(ii), because a unit in~$\princ{\ell}{n}$ then acts trivially on an element killed
by~$\ell^{k_0}$. And every element of $V$ is the product of a power of~$\eta$ and an element
of~$\levU{\ell}{H}{n}$. For \m{H=\GalG{F}{\Omega}} the group $\levU{\ell}{H}{n}$ is the group
$\Un{F}{\Omega}{\ell}{n}$ of Definition~\ref{d:groups}.

% ------------------------------------------------------------------------------------- 6
\sect{Notes on the formal text}{s:notes}
\begin{enumerate}[label=\textup{\arabic*.},leftmargin=1.6em,itemsep=5pt]
\item \emph{The conditions at the places are not used.} The formal statement takes $\varphi$
in~$\Dinf{\ell}{N}{H}$. Its
proof unfolds this hypothesis once and keeps two of its four clauses: $\varphi$ is locally constant, and it is
additive. The clauses on inertia and decomposition groups are not referred to again, and neither is the
integer~$N$. So Theorem~A holds for every locally constant homomorphism \m{\cycU{\ell}{H}\to\QZl{\ell}}. The
conditions matter for Proposition~\ref{p:action}(a), which makes $\Dinf{\ell}{N}{H}$ a module with an action
of~$\GalG{F}{\Omega}$.
\item \emph{The exponent.} The formal statement asserts, for each~$\gamma$, that some~$k$ has the two properties
of~(ii). The proof constructs \m{k=\ell^nk_0}, the same for all~$\gamma$; this is the form stated above.
\item \emph{Where the hypotheses on $H$ enter.} Closedness gives the compactness of~$\cycU{\ell}{H}$, used for the
open subgroup~$W$ and for the finitely many values, and the compactness of the set~$K$. Normality makes
$\cycU{\ell}{H}$ normal, so that conjugation acts.
\item \emph{The coefficients.} The module $\QZl{\ell}$ enters through Lemma~\ref{l:four}(c) alone. The argument of
§\ref{s:proof} applies to locally constant homomorphisms into any module whose elements are killed by powers
of~$\ell$. This remark is ours; the archive states the theorem for~$\QZl{\ell}$.
\item \emph{The polynomial ideal} is the device of the formal proof, followed here step by step. The sketch in the
first note used the same congruence, applied $k_0$ times to elements; the ideal does this bookkeeping once.
\item \emph{The rest of the unit.} The unit of the archive that holds these declarations has \UnitMembers{}
declarations that the final theorems use. This note treats \UnitInNote. The others are, by source file:

\par\vspace{2pt}
{\fontsize{9.2}{12}\selectfont\noindent\hspace*{1.2em}\begin{tabular}{@{}l@{\hspace{14pt}}r@{}}
\ElsewhereRows
\end{tabular}\par}
\vspace{2pt}

The grouping is provisional. Each group is the subject of a note to come.
\end{enumerate}

\subhead{How this note was written.} The reading project supplied a dossier of the unit: the statements, a
glossary, and a structured proof of every theorem in which each step quotes the lines of the archive that it
renders. A program compared the quotations with the archive, for all \UnitSteps{} steps of the \UnitProofs{}
proofs of the unit, and found that no line of these proofs is left without a step. The proofs of this note were
written from those structured proofs. The Lean files of the unit were not opened then; for this edition the
definitions were read there, and the proofs were not (see below). The statement of item~1 rests
on the quoted lines of the structured proof of the theorem: the hypothesis on~$\varphi$ occurs in four of them
and nowhere else. The statements quoted below were read from the archive by the build script, and they agree
with the dossier's.

\subhead{Changes in this edition.} The first edition fixed the prime, the number field and its algebraic closure
once, at the head of §\ref{s:setting}, and then wrote the cyclotomic character with one letter, the Galois group
with one letter, the groups of the tower with a single index, the module of characters without its ring of
coefficients, the two modules at the top of the tower without any argument, and the action by conjugation as a
product. That hid arguments of the formal definitions: the fields and the prime in the character and in the
groups, the ring in the module of characters, the prime, the integer~$N$ and the subgroup in the modules at the
top of the tower, and the representation in which an element acts. In this edition every function that the
archive defines is written with all the explicit arguments of its formal definition, in every place; the
arguments were read from the Lean text of each declaration. The definitions that stood in running text are
numbered definitions now, so the items of §\ref{s:setting} are renumbered, and every theorem names the prime, the
field and its closure in its own statement. Some things are stated that were only described before: the two
modules of Definition~\ref{d:equiv}, through which Proposition~\ref{p:action}(c) is phrased as the archive
phrases it; and in §\ref{s:use} the two uses of the theorem, with the proposition of the first note that the first
of them needs. The second use is quoted without proof. Three letters were renamed, in Definition~\ref{d:char} and
in the proof of Lemma~\ref{l:four}, so that none has two meanings. No statement and no proof of the first edition
changed in content.

% ------------------------------------------------------------------------------------- A
\head{Appendix A. Concordance with the Lean archive}
{\fontsize{9.6}{12.8}\selectfont\color{inksoft}\noindent Each name is a link to the lines of the declaration at
commit {\leanfont\Pin} of the public mirror of the archive. The line ranges and the quotations below are read from
the archive by the build script, which also checks them against the SHA-256 of each file. The first thirteen rows
are the declarations of the unit that this note treats.\par}
\vspace{4pt}
{\fontsize{8.6}{11.6}\selectfont\renewcommand{\leanfont}{\ttfamily\fontsize{7.7}{10}\selectfont\color{leancol}}
\setlength{\LTleft}{0pt}\setlength{\LTright}{\fill}\setlength{\tabcolsep}{0pt}
\begin{longtable}{@{}p{30mm}p{56mm}p{43mm}r@{}}
\textbf{In this note} & \textbf{Lean declaration} & \textbf{File} & \textbf{Lines}\\\midrule\endhead
\ConcordanceRows
\end{longtable}}

\Needspace{8\baselineskip}
\subhead{Three statements as they are written in the archive.} Theorem~A; here {\leanfont U} is the group~$H$,
{\leanfont levU l U n} is~$\levU{\ell}{H}{n}$, {\leanfont cycU l U} is~$\cycU{\ell}{H}$, and {\leanfont funRep}
with its three arguments is the representation~$\pi$:
\quoted{dinfNil}
The module of characters \m{\XDual{N}{C}{R}{M}}, where {\leanfont U} is the group~$C$:
\quoted{XDual}
Proposition~\ref{p:action}:
\quoted{funRepMem}

% ------------------------------------------------------------------------------------- notation
\head{Notation}
{\fontsize{9.6}{12.6}\selectfont\color{inksoft}\noindent The symbols of the text. A name in brackets is the one
used in the archive, with the explicit arguments of its definition there.\par}
\vspace{4pt}
\begingroup\fontsize{8.6}{10.9}\selectfont\raggedright
\renewcommand{\leanfont}{\ttfamily\fontsize{7.6}{10}\selectfont\color{leancol}}
\setlength{\columnsep}{9mm}\setlength{\parskip}{0pt}
\begin{multicols}{2}
\nt{$\ell$, $\Z_\ell$, $\Q_\ell$}{a prime number; the ring of $\ell$-adic integers; its field of fractions}
\nt{$\QZl{\ell}$}{the quotient group of $\Q_\ell$ by~$\Z_\ell$ [\mbox{\leanfont QZ l}]}
\nt{$F$, $\Omega$, $E$, $F_\infty$}{a number field; an algebraic closure of~$F$; a finite extension of~$F$; $F$ with all roots of unity of $\ell$-power order}
\nt{$\GalG{F}{\Omega}$}{the Galois group of $\Omega$ over~$F$ [\mbox{\leanfont G F Ω}]}
\nt{$\cyc{F}{\Omega}{\ell}{g}$}{the value at~$g$ of the $\ell$-adic cyclotomic character of $\Omega$ over~$F$ [\mbox{\leanfont cycZ F Ω l g}]; with a set in the last place, the set of the values on it}
\nt{$\zeta$}{a root of unity}
\nt{$\princ{\ell}{n}$}{the units \m{1+\ell^n\Z_\ell}, with \m{\princ{\ell}{0}=\Z_\ell^\times} [\mbox{\leanfont princ l a}]}
\nt{$H$, $\levU{\ell}{H}{n}$, $\cycU{\ell}{H}$}{a subgroup of~$\GalG{F}{\Omega}$; its elements whose character value lies in~$\princ{\ell}{n}$ [\mbox{\leanfont levU l U n}]; its elements whose character value is~$1$ [\mbox{\leanfont cycU l U}]}
\nt{$\Un{F}{\Omega}{\ell}{n}$}{the group $\levU{\ell}{H}{n}$ for the whole group \m{H=\GalG{F}{\Omega}} [\mbox{\leanfont Un F Ω l n}]}
\nt{$C$, $N$, $R$, $M$}{a subgroup of~$\GalG{F}{\Omega}$; an integer \m{\ge0}; a commutative ring; an $R$-module}
\nt{$A$, $\mathfrak m_A$, $z$}{a valuation ring of~$\Omega$; its maximal ideal; an element of~$\Omega$}
\nt{$\Dec{A}$, $\Ine{A}$}{the decomposition group and the inertia group of~$A$ in~$\GalG{F}{\Omega}$ [\mbox{\leanfont Dec A}, \mbox{\leanfont Ine A}]}
\nt{$\XDual{N}{C}{R}{M}$}{the characters of~$C$ with values in~$M$ (Definition~\ref{d:char}) [\mbox{\leanfont XDual N U R M}]}
\nt{$\Dinf{\ell}{N}{H}$, $\Xinf{\ell}{N}{H}$}{the module \m{\XDual{N}{\cycU{\ell}{H}}{\Z_\ell}{\QZl{\ell}}} [\mbox{\leanfont Dinf l N U}]; its dual [\mbox{\leanfont Xinf l N U}]}
\nt{$\funRep{R}{M}{C}$}{the functions \m{C\to M} with the action of~$\GalG{F}{\Omega}$ by conjugation (Definition~\ref{d:conj}) [\mbox{\leanfont funRep R M U}]}
\nt{$\pi$}{the representation \m{\funRep{\Z_\ell}{\QZl{\ell}}{\cycU{\ell}{H}}} (Theorem~A)}
\nt{$\conjCond{C}{\sigma}{\theta}$, $\EqHom{N}{C}{\rho}$}{modules of functions \m{C\to M} (Definition~\ref{d:equiv}) [\mbox{\leanfont conjCond U σ T}, \mbox{\leanfont EqHom N U ρ}]}
\nt{$\sigma$, $\theta$, $\rho$}{an automorphism of the field~$\Omega$; an $R$-linear map \m{M\to M}; a representation of~$\GalG{F}{\Omega}$ on~$M$}
\nt{$\varphi$}{a character, or any function on a subgroup}
\nt{$g$, $h$, $\gamma$, $\eta$, $x$, $y$, $c$, $w$, $a$, $b$}{elements of~$\GalG{F}{\Omega}$}
\nt{$u$, $v$, $s$, $d$, $r$}{$\ell$-adic numbers; \m{s=1+v+\dots+v^{\ell-1}}}
\nt{$W$, $S$, $K$, $V$}{an open normal subgroup; the set $WC$; the set \m{H\smallsetminus S}; a subgroup of~$\GalG{F}{\Omega}$ (§\ref{s:use})}
\nt{$n$, $m$, $j$, $i$, $k$, $k_0$, $q$, $e$}{integers; \m{\ell^{k_0}\varphi=0}; \m{q=\ell^n}; $e$ is the exponent of the twist in Proposition~\ref{p:inv}}
\nt{$T$, $J$, $P$, $Q$, $Y$}{the operator \m{\pi(\gamma)-1}; the ideal of polynomials that kill~$\varphi$; polynomials; the indeterminate}
\nt{$B$, $f$, $t$, $\xi$, $\beta$}{a module over~$\Z_\ell$; an endomorphism of~$B$; an endomorphism of the dual \m{\Hom_{\Z_\ell}(B,\QZl{\ell})}; an element of the dual; an element of~$B$}
\nt{$L$}{the number field of the first note}
\end{multicols}
\endgroup

% ------------------------------------------------------------------------------------- references
\head{References}
\begingroup\fontsize{9.4}{12.6}\selectfont\raggedright\renewcommand{\section}[2]{}
\begin{thebibliography}{NSW}\setlength{\itemsep}{2pt}
\bibitem[Al]{Al} L.\,Alpöge, announcement of proofs of the Riemann hypothesis and of the generalized Riemann
  hypothesis for Dirichlet characters, 20 November 2026, with a Lean archive;
  \href{\PinUrl}{public mirror of the archive at commit {\leanfont\Pin}}.
\bibitem[N1]{N1} \emph{Twisted cokernels in a cyclotomic tower}, Deformalization~1, riemann-rev,
  \href{https://nasqret.github.io/riemann-rev/digests/tower-bound.pdf}{\texttt{digests/tower-bound.pdf}}.
\bibitem[NSW]{NSW} J.\,Neukirch, A.\,Schmidt and K.\,Wingberg, \emph{Cohomology of Number Fields}, 2nd ed.,
  Grundlehren der mathematischen Wissenschaften 323, Springer, 2008.
\end{thebibliography}
\endgroup

\end{document}
